\documentclass[10pt,a4paper]{amsart}
\usepackage{amsmath,amssymb,amsthm,graphicx}
\usepackage{verbatim}
\newenvironment{reference}{\comment}{\endcomment}
%\newenvironment{reference}{}{}
\newenvironment{slogan}{\comment}{\endcomment}
\newenvironment{history}{\comment}{\endcomment}

% For commutative diagrams we use Xy-pic
\usepackage[all]{xy}

% We use 2cell for 2-commutative diagrams.
\xyoption{2cell}
% Theorem environments.
%
\theoremstyle{plain}
\newtheorem{theorem}[subsection]{Theorem}
\newtheorem{proposition}[subsection]{Proposition}
\newtheorem{lemma}[subsection]{Lemma}

\theoremstyle{definition}
\newtheorem{definition}[subsection]{Definition}
\newtheorem{example}[subsection]{Example}
\newtheorem{exercise}[subsection]{Exercise}
\newtheorem{situation}[subsection]{Situation}

\theoremstyle{remark}
\newtheorem{remark}[subsection]{Remark}
\newtheorem{remarks}[subsection]{Remarks}

\numberwithin{equation}{subsection}

% Macros
%
\def\lim{\mathop{\mathrm{lim}}\nolimits}
\def\colim{\mathop{\mathrm{colim}}\nolimits}
\def\Spec{\mathop{\mathrm{Spec}}}
\def\Hom{\mathop{\mathrm{Hom}}\nolimits}
\def\Ext{\mathop{\mathrm{Ext}}\nolimits}
\def\SheafHom{\mathop{\mathcal{H}\!\mathit{om}}\nolimits}
\def\SheafExt{\mathop{\mathcal{E}\!\mathit{xt}}\nolimits}
\def\Sch{\mathit{Sch}}
\def\Mor{\mathop{\mathrm{Mor}}\nolimits}
\def\Ob{\mathop{\mathrm{Ob}}\nolimits}
\def\Sh{\mathop{\mathit{Sh}}\nolimits}
\def\NL{\mathop{N\!L}\nolimits}
\def\CH{\mathop{\mathrm{CH}}\nolimits}
\def\proetale{{pro\text{-}\acute{e}tale}}
\def\etale{{\acute{e}tale}}
\def\QCoh{\mathit{QCoh}}
\def\Ker{\mathop{\mathrm{Ker}}}
\def\Im{\mathop{\mathrm{Im}}}
\def\Coker{\mathop{\mathrm{Coker}}}
\def\Coim{\mathop{\mathrm{Coim}}}

% Boxtimes
%
\DeclareMathSymbol{\boxtimes}{\mathbin}{AMSa}{"02}

%
% Macros for moduli stacks/spaces
%
\def\QCohstack{\mathcal{QC}\!\mathit{oh}}
\def\Cohstack{\mathcal{C}\!\mathit{oh}}
\def\Spacesstack{\mathcal{S}\!\mathit{paces}}
\def\Quotfunctor{\mathrm{Quot}}
\def\Hilbfunctor{\mathrm{Hilb}}
\def\Curvesstack{\mathcal{C}\!\mathit{urves}}
\def\Polarizedstack{\mathcal{P}\!\mathit{olarized}}
\def\Complexesstack{\mathcal{C}\!\mathit{omplexes}}
% \Pic is the operator that assigns to X its picard group, usage \Pic(X)
% \Picardstack_{X/B} denotes the Picard stack of X over B
% \Picardfunctor_{X/B} denotes the Picard functor of X over B
\def\Pic{\mathop{\mathrm{Pic}}\nolimits}
\def\Picardstack{\mathcal{P}\!\mathit{ic}}
\def\Picardfunctor{\mathrm{Pic}}
\def\Deformationcategory{\mathcal{D}\!\mathit{ef}}
\usepackage{cleveref}
\crefname{originaldefinition}{original definition}{original definitions}
\Crefname{originaldefinition}{Original definition}{Original definitions}
\makeatletter
\newlabel{categories-definition-system-over-poset}{{0030}{source}{}{sourcepointer-1}{}}
\newlabel{categories-definition-system-over-poset@cref}{{[originaldefinition][0][]0030}{[1][0][]source}{}{}{}}
\makeatother
\crefname{originalsection}{original section}{original sections}
\Crefname{originalsection}{Original section}{Original sections}
\makeatletter
\newlabel{section-points}{{00Y3}{source}{}{sourcepointer-2}{}}
\newlabel{section-points@cref}{{[originalsection][0][]00Y3}{[1][0][]source}{}{}{}}
\makeatother
\crefname{originallemma}{original lemma}{original lemmas}
\Crefname{originallemma}{Original lemma}{Original lemmas}
\makeatletter
\newlabel{categories-lemma-directed-commutes}{{002W}{source}{}{sourcepointer-3}{}}
\newlabel{categories-lemma-directed-commutes@cref}{{[originallemma][0][]002W}{[1][0][]source}{}{}{}}
\makeatother
\crefname{originallemma}{original lemma}{original lemmas}
\Crefname{originallemma}{Original lemma}{Original lemmas}
\makeatletter
\newlabel{lemma-maps-u-points}{{00YH}{source}{}{sourcepointer-4}{}}
\newlabel{lemma-maps-u-points@cref}{{[originallemma][0][]00YH}{[1][0][]source}{}{}{}}
\makeatother
\crefname{originalsection}{original section}{original sections}
\Crefname{originalsection}{Original section}{Original sections}
\makeatletter
\newlabel{categories-section-directed-colimits}{{04AX}{source}{}{sourcepointer-5}{}}
\newlabel{categories-section-directed-colimits@cref}{{[originalsection][0][]04AX}{[1][0][]source}{}{}{}}
\makeatother
\crefname{originallemma}{original lemma}{original lemmas}
\Crefname{originallemma}{Original lemma}{Original lemmas}
\makeatletter
\newlabel{lemma-refine}{{00YO}{source}{}{sourcepointer-6}{}}
\newlabel{lemma-refine@cref}{{[originallemma][0][]00YO}{[1][0][]source}{}{}{}}
\makeatother
\title{Refining a system at once for all finite coverings}
\author{The Stacks Project Authors}
\date{stacks-project:sites.tex:local-snapshot-20261012; GNU FDL 1.2 or later}
\begin{document}
\maketitle
\noindent\textit{Editorial note.} All delivered mathematical text is the author's own, copied from the pinned Stacks Project chapter file sites-cohomology.tex; no statement or proof step was rewritten, reordered or removed.\par\medskip
\noindent\textit{Editorial note.} The proof is the author's complete argument; the informal Zorn step is immediately replaced by the explicit transfinite recursion, and the remaining verification is the routine check that the union has the required properties, which is stated.\par\medskip
\noindent\textit{Editorial note.} Every labelled reference occurring inside the counted spans is carried in assembly.external\_references as a non-load-bearing navigation pointer (same-chapter ones at their definition site, cross-chapter ones at the invocation inside the counted span), so no load-bearing mathematics is replaced by a stub. No bibliography citation occurs inside the counted spans.\par\medskip
\noindent\textit{Editorial note.} The checked chapter list for same-name collisions is the set of stacks chapters staged in this repository: chow, cohomology, divisors, etale-cohomology, flat, more-algebra, more-morphisms, morphisms, spaces-more-morphisms, varieties together with this card's chapter sites-cohomology.\par\medskip
\noindent\textit{Editorial note.} The preamble is the chapter's own retained minimal preamble (sources/stacks-sites-cohomology/retained-preamble.tex), a verbatim excerpt of the native preamble: the Xy-pic setup and the native theorem declarations with the mathematical macros.\par\medskip
\noindent\textit{Allowed standard prerequisites (editorial).}\begin{itemize}
\item One-step refinement of a system. Same chapter, Lemma lemma-refine. Supplies the single step of the transfinite construction.
\item Directed colimits of sets. Categories, Section categories-section-directed-colimits. Supplies the description of the values of the functor as directed colimits and the distinctness of the images of the two sections.
\end{itemize}\medskip
\ifdefined\hypertarget\hypertarget{sourcepointer-1}{}\fi
\noindent\textit{Source reference.} Original definition 0030: sources/stacks-sites/pinned-original-source.tex, source lines 9466--9466. Navigation pointer invoked at proof line 9466; the counted span contains the pointer only and the cited result is not reproduced inside it, so the pointer is registered as a non-load-bearing navigation stub in the established form for Stacks material. No bibliography citation is involved.\par\medskip
\ifdefined\hypertarget\hypertarget{sourcepointer-2}{}\fi
\noindent\textit{Source reference.} Original section 00Y3: sources/stacks-sites/pinned-original-source.tex, source lines 9473--9473. Navigation pointer invoked at proof line 9473; the counted span contains the pointer only and the cited result is not reproduced inside it, so the pointer is registered as a non-load-bearing navigation stub in the established form for Stacks material. No bibliography citation is involved.\par\medskip
\ifdefined\hypertarget\hypertarget{sourcepointer-3}{}\fi
\noindent\textit{Source reference.} Original lemma 002W: sources/stacks-sites/pinned-original-source.tex, source lines 9482--9482. Navigation pointer invoked at proof line 9482; the counted span contains the pointer only and the cited result is not reproduced inside it, so the pointer is registered as a non-load-bearing navigation stub in the established form for Stacks material. No bibliography citation is involved.\par\medskip
\ifdefined\hypertarget\hypertarget{sourcepointer-4}{}\fi
\noindent\textit{Source reference.} Original lemma 00YH: sources/stacks-sites/pinned-original-source.tex, source lines 9500--9500. Navigation pointer invoked at proof line 9500; the counted span contains the pointer only and the cited result is not reproduced inside it, so the pointer is registered as a non-load-bearing navigation stub in the established form for Stacks material. No bibliography citation is involved.\par\medskip
\ifdefined\hypertarget\hypertarget{sourcepointer-5}{}\fi
\noindent\textit{Source reference.} Original section 04AX: sources/stacks-sites/pinned-original-source.tex, source lines 9589--9589. Navigation pointer invoked at proof line 9589; the counted span contains the pointer only and the cited result is not reproduced inside it, so the pointer is registered as a non-load-bearing navigation stub in the established form for Stacks material. No bibliography citation is involved.\par\medskip
\ifdefined\hypertarget\hypertarget{sourcepointer-6}{}\fi
\noindent\textit{Source reference.} Original lemma 00YO: sources/stacks-sites/pinned-original-source.tex, source lines 9596--9596. Navigation pointer invoked at proof line 9596; the counted span contains the pointer only and the cited result is not reproduced inside it, so the pointer is registered as a non-load-bearing navigation stub in the established form for Stacks material. No bibliography citation is involved.\par\medskip
\medskip\noindent
Let $\mathcal{C}$ be a site.
Suppose that $(I, \geq)$ is a directed set,
and that $(U_i, f_{ii'})$ is an inverse system over $I$, see
Categories, Definition \ref{categories-definition-system-over-poset}.
Given the data $(I, \geq, U_i, f_{ii'})$ we define
$$
u : \mathcal{C} \longrightarrow \textit{Sets}, \quad
u(V) = \colim_i \Mor_\mathcal{C}(U_i , V)
$$
Let $\mathcal{F} \mapsto \mathcal{F}_p$ be the stalk functor
associated to $u$ as in Section \ref{section-points}.
It is direct from the definition that actually
$$
\mathcal{F}_p = \colim_i \mathcal{F}(U_i)
$$
in this special case.
Note that $u$ commutes with all finite limits (I mean those that
are representable in $\mathcal{C}$) because
each of the functors $V \mapsto \Mor_\mathcal{C}(U_i , V)$
do, see Categories, Lemma \ref{categories-lemma-directed-commutes}.

\medskip\noindent
We say that a system $(I, \geq, U_i, f_{ii'})$
is a {\it refinement} of $(J, \geq, V_j, g_{jj'})$ if
$J \subset I$, the ordering on $J$ induced from that of $I$
and $V_j = U_j$, $g_{jj'} = f_{jj'}$ (in words, the inverse system
over $J$ is induced by that over $I$). Let $u$ be the functor
associated to $(I, \geq, U_i, f_{ii'})$ and let $u'$ be the
functor associated to $(J, \geq, V_j, g_{jj'})$.
This induces a transformation of functors
$$
u' \longrightarrow u
$$
simply because the colimits for $u'$ are over a subsystem
of the systems in the colimits for $u$.
In particular we get an associated transformation of
stalk functors $\mathcal{F}_{p'} \to \mathcal{F}_p$,
see Lemma \ref{lemma-maps-u-points}.

\begin{lemma}
\label{lemma-refine-all-at-once}
Let $\mathcal{C}$ be a site.
Let $(J, \geq, V_j, g_{jj'})$ be a system as above with associated
pair of functors $(u', p')$.
Let $\mathcal{F}$ be a sheaf on $\mathcal{C}$.
Let $s, s' \in \mathcal{F}_{p'}$ be distinct elements.
There exists a refinement $(I, \geq, U_i, f_{ii'})$
of $(J, \geq, V_j, g_{jj'})$ such that $s, s'$ map
to distinct elements of $\mathcal{F}_p$ and such that
for every finite covering $\{W_k \to W\}$ of the site
$\mathcal{C}$, and any $f \in u'(W)$ the image of $f$ in $u(W)$
is in the image of one of the $u(W_k)$.
\end{lemma}

\begin{proof}
Let $E$ be the set of pairs $(\{W_k \to W\}, f\in u'(W))$.
Consider pairs $(E' \subset E, (I, \geq, U_i, f_{ii'}))$
such that
\begin{enumerate}
\item $(I, \geq, U_i, g_{ii'})$ is a refinement of $(J, \geq, V_j, g_{jj'})$,
\item $s, s'$ map to distinct elements of $\mathcal{F}_p$, and
\item for every pair $(\{W_k \to W\}, f\in u'(W)) \in E'$ we have that
the image of $f$ in $u(W)$ is in the image of one of the $u(W_k)$.
\end{enumerate}
We order such pairs by inclusion in the first factor and
by refinement in the second. Denote $\mathcal{S}$ the class
of all pairs $(E' \subset E, (I, \geq, U_i, f_{ii'}))$ as above.
We claim that the hypothesis of Zorn's lemma holds for $\mathcal{S}$. Namely,
suppose that $(E'_a, (I_a, \geq, U_i, f_{ii'}))_{a \in A}$
is a totally ordered subset of $\mathcal{S}$. Then we can define
$E' = \bigcup_{a \in A} E'_a$ and we can set $I = \bigcup_{a \in A} I_a$.
We claim that the corresponding pair
$(E' , (I, \geq, U_i, f_{ii'}))$ is an element of $\mathcal{S}$.
Conditions (1) and (3) are clear. For condition (2) you note
that
$$
u = \colim_{a \in A} u_a
\text{ and correspondingly }
\mathcal{F}_p = \colim_{a \in A} \mathcal{F}_{p_a}
$$
The distinctness of the images of $s, s'$ in this stalk follows
from the description of a directed colimit of sets, see
Categories, Section \ref{categories-section-directed-colimits}.
We will simply write
$(E', (I, \ldots)) = \bigcup_{a \in A}(E'_a, (I_a, \ldots))$
in this situation.

\medskip\noindent
OK, so Zorn's Lemma would apply if $\mathcal{S}$ was a set,
and this would, combined with Lemma \ref{lemma-refine} above easily prove
the lemma. It doesn't since $\mathcal{S}$ is a class. In order
to circumvent this we choose a well ordering on $E$.
For $e \in E$ set $E'_e = \{e' \in E \mid e' \leq e\}$.
Using transfinite recursion we construct pairs
$(E'_e, (I_e, \ldots)) \in \mathcal{S}$ such that
$e_1 \leq e_2 \Rightarrow (E'_{e_1}, (I_{e_1}, \ldots))
\leq (E'_{e_2}, (I_{e_2}, \ldots))$.
Let $e \in E$, say $e = (\{W_k \to W\}, f\in u'(W))$.
If $e$ has a predecessor $e - 1$, then we let
$(I_e, \ldots)$ be a refinement of $(I_{e - 1}, \ldots)$
as in Lemma \ref{lemma-refine} with respect to the system
$e = (\{W_k \to W\}, f\in u'(W))$.
If $e$ does not have a predecessor, then we let
$(I_e, \ldots)$ be a refinement of $\bigcup_{e' < e} (I_{e'}, \ldots)$
with respect to the system
$e = (\{W_k \to W\}, f\in u'(W))$.
Finally, the union $\bigcup_{e \in E} I_e$ will be a solution to
the problem posed in the lemma.
\end{proof}

\end{document}
